\chapter{取り込まなかった資料と、会話で検討したこと}\label{app:others}

この付録は 2 つの部分からなる。前半は、資料集のうちこの本に取り込まなかったファイルと、その中身がこの本のどこにあるかの対応である。後半は、文書にはしなかったが、作業の途中の会話で検討し、方針に影響した事柄の記録である。

\section{取り込まなかった資料}

\begin{longtable}{P{5.6cm}P{8.8cm}}
\toprule
\textbf{ファイル} & \textbf{中身と、この本での場所} \\
\midrule
\endhead
\bottomrule
\endfoot
\code{04\_設計書/LLM-BLM\_DB設計.html} & 最初の DB の設計（10/13 まで）。表と列の定義は、付録「テーブルとビューの定義」（総合報告書 1 の付録）と同じ内容で、その後の拡張は第 VI 部にある \\
\code{06\_中間発表と途中のメモ/} の HTML 12 本 & 総合報告書 1 を書く前の途中のメモ（テスト実行の結果、データの扱い方、中間発表の流れ、予測 LLM の入出力の設計、今後の計画、元論文の計算の問題点、実験 F、実験のまとめ、段階 1 の結果、解説、評価ポイント）。どれも総合報告書 1 の同じ名前の章にまとめ直されている（第 VII 部） \\
\code{06\_中間発表と途中のメモ/} の PPTX & 中間発表のスライド（10/13 まで）。内容は総合報告書 1 の前半と同じ \\
\code{06\_中間発表と途中のメモ/LLM-BLM\_データ確認.ipynb} & データの確認のノートブック。結果は第 VI 部のデータ基盤の章にある \\
\code{07\_参考文献/} の PDF 3 本 & 元論文（arXiv:2504.14345 の第 2 版）、霧生・枇々木（2014）、ARGUS（2026）。他人の著作物なので取り込まない。後の 2 本は、この付録の後半で中身を説明する \\
\code{03\_方針と文献調査/今後の計画/渡した資料/} & 計画の依頼に添付した、総合報告書 2 の原稿と \code{やること.md} の写し。元の文書を取り込んだので省いた \\
\code{03\_方針と文献調査/読ませ方とまとめ方/deep/文献照合.md} & 文献の照合のスクリプトが、arXiv の API の時間切れで止まったときの出力。照合は手で行い、結果は比較検討の文書にある（第 X 部） \\
\code{experiments/*/results/} の CSV・JSON & 実験の数値の元のデータ。表と図に使った値は、各実験の章にある \\
\end{longtable}

\section{会話で検討したこと}

\subsection{オプションから分布を読む：霧生・枇々木（2014）}

\begin{paperbox}[霧生拓也・枇々木規雄「複数資産にインプライド分布を用いた最適資産配分モデル」（日本オペレーションズ・リサーチ学会和文論文誌 57 巻、2014）]
オプションの価格から逆算できる原資産の確率分布（\term{インプライド分布}）を資産配分に使う研究である。Black--Scholes 型の価格式を逆に解くと、市場が織り込んでいる将来の価格の分布（リスク中立の分布）が出てくる。過去のデータから推定する分布（ヒストリカル分布）と違い、判断の時点の市場の見通しが入っている。論文は、(1) 過去のデータでリスク調整（リスク中立の分布を現実の分布に直すこと）をすると、かえって成績が下がる、(2) 満期の違いを薄板スプラインで平滑化して取り除き、複数の資産にインプライド分布を使うと成績が上がる、と報告している。
\end{paperbox}

この論文を読んだあと、「LLM の見通しとオプションの見通しを、両方ベイズで混ぜればよいのでは」という案が出た。個別の銘柄やセクターのオプションの過去の価格は無料では取れない（総合報告書 1 のセクターの指標の章）ので、市場全体のオプションから出る VIX だけを使い、共分散の大きさを VIX に合わせる補足実験 X04 を行った。結果は、TE の校正（X03）に足しても改善しなかった（2020 年 2 月の急落の月だけ効いた）。

\subsection{ARGUS（2026）}

\begin{paperbox}[Shu・Lu・Zhang・Sun「ARGUS: Augmented Reasoning Over Multi-Modal Text and Transaction Streams for Robust Black--Litterman Portfolio Optimization」（IEEE Access、2026）]
BLM の 3 つの入力（期待リターン、共分散、運用の制約）を、LLM で同時に作る枠組みである。
\begin{itemize}
\item \textbf{ビューの生成}：文章の証拠と、取引の細かいデータ（売買の向き）を突き合わせ、食い違うときはビューを弱める。
\item \textbf{共分散の補強}：高頻度のデータから計算した実現の共分散に、LLM が見込む相関の変化を足す。
\item \textbf{制約の翻訳}：自然言語の運用方針（「流動性の低い銘柄は控える」など）を、数式の制約に直す。
\item \textbf{閉じた較正}：結果を見てビューの確からしさを直す。複数の情報源のビューを、外れ値に強い幾何中央値でまとめる。
\end{itemize}
\SP{} のセクター ETF、大型株、暗号資産で、7 つの比較対象を上回ったと報告している。
\end{paperbox}

報告されている成績は、この研究では確かめていない。この研究の方針に取り入れたのは考え方の 2 つで、どちらも次に確かめる候補に入れた（\chapref{ch:open}）。
\begin{itemize}
\item 「文章と売買の向きの食い違い」を、日次のデータに直した形（文章の点数と、発表日の株価の向きが食い違った発表）。
\item 「情報源どうしの意見の割れ」でビューの確からしさ（$\Omega$）を決める。
\end{itemize}

\subsection{マクロの変数は要るか}

「GDP・インフレ・金利をまったく考慮していないが、株価を見れば同じことか」という問いが出た。答えは「半分は同じ」である。
\begin{itemize}
\item 公表されたマクロの数字（とその予想）は、公表と同時に株価に入る。株価を使えば、間接的にマクロの情報を使っていることになる。BLM の均衡リターンも、時価総額（市場の評価）から出すので、市場がマクロをどう見ているかを含む。
\item ただし、マクロの変数で「どの業種が得をするか」の関係は、銘柄の株価だけからは学びにくい。セクター単位の見通しには使える可能性がある（第 IV 部「マクロ経済とボラティリティ」）。
\end{itemize}
追加のデータとして FRED の 17 系列を DB に入れたが、実験には使っていない。

\subsection{データの期間と上場廃止の銘柄}

「30 年前の株価と、上場をやめた銘柄の株価は取れないか」を調べた。yfinance には上場を廃止した銘柄の株価がないことが多い。EODHD では、2015 年以降に指数にいて yfinance に株価のない 157 銘柄のうち 134 銘柄が上場廃止の一覧にあったが、取れる期間は約 12 年分までだった。大学の契約のデータベースも見つからなかった。結論として、データは今のまま（2015 年以降、各判断日の構成銘柄、株価のない銘柄は除く）とし、生存者バイアスを限界として書くことにした（\chapref{ch:pit}）。

\subsection{2015 年より前の文章で学習できるか}

決算発表文やニュースは 2015 年より前にもある。ただし、それで「反応」を学習するには、同じ時期の株価（正解）が要る。この研究の DB の株価は 2015 年からなので、それより前の文章は反応の学習には使えない。正解のない文章で学ぶこと（事前学習）は、ChronoGPT がすでに行っている。

\subsection{人工データ、LLM の追加学習、計算の場所}

「人工データで水増しすれば深層学習ができるのでは」「LLM 自体をファインチューニングできないか（PEFT、ソフトプロンプト）」「内蔵の GPU も使えば速くなるか」「量子化しなくてよかったのでは」という問いが出た。答えはそれぞれ\chapref{ch:ml}と\chapref{ch:llm}に書いた。LLM の追加学習は、実験 10（LoRA）として行った。

\subsection{アクティブ運用の考え方}

Pedersen『Efficiently Inefficient』（邦訳『ヘッジファンドのアクティブ投資戦略』）について、「偉人の言葉と数式が混ざっていて読みにくい」「期待リターンもリスクプレミアムの話ばかりで、市場が落とした金を拾うのがアクティブ運用の本質か」という話をした。この本の立場は、市場は「効率的に非効率」で、情報を集めて処理する費用に見合うだけの小さな非効率が残り、それを拾う者が報われる、というものである。ファクター（割安、モメンタム、低リスクなど）のリスクプレミアムを、どの期間で、どの費用で拾うかが運用の中身になる。この研究の結論（LLM の読み取りは 1 日で価格に入る）は、文章の情報の処理の費用が、LLM によって大きく下がったことの裏返しとも読める。期間を長くとる（月次・四半期）分析が向く理由も、ここにある（\chapref{ch:active}）。
